% ======================================================================
% 第 3 章 热分析方法  分片 b
% 源：sources/parts/ch03-b.txt（上册 PDF p.200--237 / 书内 p.181--218）
% 边界：分片 a 已写热分析分类、TG/DTA/DSC 原理与动力学简介（Ozawa-Flynn-Wall、
%       Kissinger、Freeman-Carroll 三种处理方法的推导）。本分片不重复，
%       直接进入应用、联用技术、实验条件与误判。
% ======================================================================

\subsection{联用热分析：质量信号必须与逸出物信号对接}

\begin{probout}
热重单独给出的是一条质量随温度（或时间）连续变化的曲线，它能回答的问题是
「在什么温度区间、失去了或增加了多少质量」，不能回答「逸出的是什么物种」。
把 TG 的失重直接翻译成某个具体化合物，是本方法最常被跨过去的一步。

联用热分析（TG-MS、TG-FTIR、TG-DSC-MS、TG-GC/MS）把质量信号与逸出物的
质荷比或红外吸收放到同一条时间轴上，输出三类信息：
\begin{itemize}
  \item 逸出物的身份，或至少是官能团类别；
  \item 每一步失重与逸出物种之间的对应关系；
  \item 由质量平衡出发的定量归属，即这一步失重折合成多少摩尔某物种。
\end{itemize}
只有第三类信息成立，TG 才从「现象记录」变成「组成与机理证据」。
\end{probout}

\begin{principle}
联用系统由三段构成：样品处的热分析主机、把逸出物从常压载气流送入高真空或
光管的传输接口、以及分析器（质谱或红外）。三段串联带来三条原理性约束，
它们共同划定了联用数据的可信区间。

\textbf{（1）传输滞后。} 逸出物在样品处生成、在分析器处被检出，中间隔着传输
线的死体积 $V$。载气流量为 $Q$ 时，平均滞后时间为
\begin{equation}
\tau \approx \frac{V}{Q},
\label{eq:ch03b-tau}
\end{equation}
若热分析以升温速率 $\beta$ 线性升温，则同一失重步在 TG 曲线与 MS 通道上的
峰位存在表观温移
\begin{equation}
\Delta T \approx \beta\,\tau .
\label{eq:ch03b-dT}
\end{equation}
这两个式子说明：联用实验里「TG 峰温」与「MS 峰温」本来就不该相等，
且升温速率越快、滞后造成的温移越大。\tocheck{未对照原页面图像：式
\equref{eq:ch03b-tau} 与 \equref{eq:ch03b-dT} 是接口死体积与传输滞后的标准
估计式，源 OCR 未给出该推导，仅在第 3 章末提到联用软件可处理 TG-MS 与
TG-FTIR 数据；此处按学科通用形式重建。}

\textbf{（2）接口温度窗。} 传输线温度必须同时满足两个相反要求：高于所有
可能逸出物的露点以免冷凝，低于最不稳定逸出物的分解温度以免二次反应。
这两条给出的窗口并不总是存在，半挥发物与热不稳定产物就落在这个缝里。

\textbf{（3）分压与稀释。} 载气把逸出物稀释几个数量级，且稀释比随流量变化。
质谱离子流强度正比于分析器处的分压而不是样品处的逸出量，所以信号强度与
失重速率之间差一个必须标定的比例因子。
\end{principle}

\begin{measure}
联用实验的必备记录项，缺一项则该数据不可定量：传输线与接口温度、载气种类
与流量、分流比（TG-DSC-MS 的分流尤其重要）、质谱扫描速率与电离方式
（EI 常用 \qty{70}{eV}）、红外光管温度与光程、以及取样触发方式
（连续采集还是定时取样）。定量时还需要同组分标准物质的响应因子标定，
即在同一接口条件下测得「单位摩尔逸出物对应的离子流积分」。
\tocheck{未对照原页面图像：本段的接口参数与标定要求为学科通用实践，
源 OCR 未列出具体数值。}
\end{measure}

\begin{dataflow}
从联用原始数据到结论，按四步走：

① 以 DTG 峰把失重过程分解成若干步，记下每步的温度区间与失重量；

② 在相同时间轴上找出每个 DTG 峰对应的质谱通道或红外吸收带；

③ 用质量平衡做交叉验算：该步实测失重（毫克）与该通道积分标定后的摩尔量
折合质量应当一致，偏差大说明有未检出物种（如高沸点物留在传输线上）或
通道重叠；

④ 只有在第 ③ 步闭合以后，才把物种名称写进结论。
\end{dataflow}

\begin{table}[tbp]
\centering
\caption{热分析联用技术的可获得信息与接口代价。读者应先看第 3 列：
接口的滞后与分流决定了第 2 列的信息能不能定量；第 4 列给出不适用的场合。
\src{本表为本改编本归纳，非源书表格；依据上册 §3.5，书内 p.201—209 / PDF p.220—228。
表中技术与接口参数为本改编本据学科通用认识整理，非源书实验数据。}}
\label{tab:ch03b-hyphenated}
\small
\begin{tabular}{@{}p{28mm}p{40mm}p{40mm}p{28mm}@{}}
\toprule
技术 & 可获得的信息 & 接口与传输滞后的代价 & 不适用的场合 \\
\midrule
TG-MS & $m/z$ 通道，适合小分子：\ce{H2O}（18）、\ce{CO2}（44）、\ce{CO}（28）、
\ce{NH3}（17）、\ce{SO2}（64）；灵敏度高、响应快 &
毛细管或分子束取样接口；传输线须保持足够高温度防冷凝；质谱真空度对载气
总流量敏感，流量不能开大 &
同质量数多来源（\ce{CO} 与 \ce{C2H4} 均为 28）时需靠碎片比区分；未标定则
只能定性 \\
TG-FTIR & 官能团信息：O--H、C=O、C--O、S=O、N--H 等吸收带，可区分异构物种 &
光管死体积大（数十至数百毫升），滞后最明显，$\tau$ 可达数十秒；光管温度
与窗片材料限制适用温区 &
易冷凝的强极性或高沸点物在管壁损失；混合物光谱叠加，须做差谱 \\
TG-DSC-MS & 同一坩埚内同时得到热流、质量与逸出物，峰位对应关系最好 &
需在 DSC 池上取气分流；分流比改变样品处实际流量与灵敏度，焓的定量基线
必须重做 &
分流后灵敏度下降，痕量逸出物可能低于检出限 \\
TG-GC/MS & 分离能力强，适合复杂混合物中同质量数物种的区分 &
采样是间歇式的，时间分辨率低（分钟量级），只能取若干温度点 &
不适合快速失重过程；单次取样量小，痕量组分易漏检 \\
\bottomrule
\end{tabular}

\vspace{1.5mm}
{\footnotesize 选型原则：需要「是不是这个物种」的判据时用 MS；需要「是哪一类
官能团」的判据时用 FTIR；需要在同一坩埚内把热效应与质量变化锁在一起时用
TG-DSC-MS；混合物复杂到 MS 无法区分时再上 GC/MS。}

\srcnote{联用软件与数据处理的说明见上册 §3.3，书内 p.181 / PDF p.200；
接口参数与选型原则为学科通用认识，源 OCR 无明文。}
\end{table}

\begin{pitfall}
\textbf{把离子流积分直接当摩尔量。} 不同 $m/z$ 通道的灵敏度相差可达数量级，
且碎片化使母体信号分流到多个通道。未做响应因子标定时，MS 只能给出
「有／无」与「先后顺序」，不能给出摩尔比。

\textbf{把 TG 峰与 MS 峰的温差不当作滞后处理。} 式 \equref{eq:ch03b-dT} 表明
这个温差随升温速率线性增长，它的存在是接口的必然结果，不是材料的两步反应。
判别方法：把升温速率降到原来的四分之一重测，若温差按比例缩小而失重步数不变，
则是滞后。

\textbf{定量外推的三条边界。} ① 没有同组分标定的通道不能定量；
② 逸出物在传输线上冷凝或二次反应时，MS 的定量结果偏低，且偏低的程度随
接口温度变化，这种误差无法用峰面积归一化消除；③ 联用数据能确定「这一步
逸出的是什么」，不能直接确定「固相中它原本以什么形式存在」——后者需要
XRD、元素分析或程序升温还原等独立证据。
\end{pitfall}

\begin{crosscheck}
联用结果必须做三类回查：与 XRD 或元素分析对照固相组成；与在线色谱对照
逸出物的摩尔量；与热重自身的失重量对照质量平衡。三者中任何一项对不上，
先怀疑接口而不是先怀疑材料。对本书涉及的应用，最常见的回查是
「TG 失重 + 理论化学计量失重」的直接比对，这一思路在催化剂前驱体组成的
确定中被完整使用过（见本章后文的实例）。
\end{crosscheck}

\subsection{实验条件对曲线的作用：峰温、峰形还是峰面积}

\begin{probout}
同一份样品，换一台热天平、换一个坩埚、换一批载气，曲线可以完全不同。
本节要回答的是一个可操作的问题：\textbf{面对一个给定的实验变量，被它改变的
究竟是峰温、峰形，还是峰面积？} 三者的性质并不相同：

\begin{itemize}
  \item \textbf{峰面积}（等价于总失重量或总焓变）是一个过程的守恒量，
  条件改变时相对稳健，是定量的首选依据；
  \item \textbf{峰温} $T_{\mathrm p}$ 是一个动力学温度，最容易被条件左右，
  它反映的是「在本次实验的时间尺度上过程\emph{走得最快}的温度」，即微商型
  曲线的峰顶；它\emph{不是}转化率走到一半的温度——对单步一级失重，峰温处的
  转化率由 $E/(R\,T_{\mathrm p})$ 决定，在 $E\approx\qty{180}{kJ.mol^{-1}}$、
  $T_{\mathrm p}\approx\qty{673}{K}$ 时 $\alpha(T_{\mathrm p})\approx0.61$，
  明显高于一半（图~\ref{fig:ch03b-rate} 与式~\equref{eq:ch03b-kissinger}）；
  \item \textbf{峰形}（宽度、对称性、肩峰）介于两者之间，它同时携带
  动力学信息与传热传质信息，是分辨重叠步骤的主要依据，也是最容易被
  误读的量。
\end{itemize}
\end{probout}

\begin{principle}
峰温与升温速率的关系不是经验规律，而是 Arrhenius 动力学的必然结果。
分片 a 已给出 Kissinger 关系，这里只引用其结果形式：
\begin{equation}
\frac{\mathrm{d}\left[\ln(\beta/T_{\mathrm p}^{2})\right]}{\mathrm{d}(1/T_{\mathrm p})}
= -\frac{E}{R},
\label{eq:ch03b-kissinger}
\end{equation}
式中 $T_{\mathrm p}$ 为微商型曲线（DTG、DTA、DSC）的峰温，$\beta$ 为升温速率，
$E$ 为活化能，$R$ 为摩尔气体常数。\src{上册 §3.3，书内 p.181 / PDF p.200}
\tocheck{未对照原页面图像：源 OCR 该式为
\texttt{d(Iln(势/TPZ\ ...\ d(l/几)=一刀尺}，严重失真；此处按 Kissinger
标准形式重建，$T_{\mathrm p}$、$\beta$、$E$、$R$ 的变量定义取自同页正文，
可靠。}

由式 \equref{eq:ch03b-kissinger} 可估出峰温对升温速率的敏感度：
\begin{equation}
\Delta T_{\mathrm p} \approx \frac{R\,T_{\mathrm p}^{2}}{E}\,\ln\frac{\beta_2}{\beta_1}.
\label{eq:ch03b-dTp}
\end{equation}
在 $T_{\mathrm p}\approx\qty{673}{K}$、$E\approx\qty{180}{kJ.mol^{-1}}$ 的量级下，
升温速率每翻一倍，峰温上移约 15 K；速率变化八倍，峰温可移动 40 K 以上。
\tocheck{未对照原页面图像：式 \equref{eq:ch03b-dTp} 为式
\equref{eq:ch03b-kissinger} 的微分近似，源 OCR 无此式，本书按标准推导补出。}

与此同时，峰面积不受 $\beta$ 影响：一个过程无论走得多快，只要始态与终态相同，
总质量变化与总焓变就相同。以纵轴取\emph{质量导数}的 DTG 曲线为例，
TG 曲线写作 $m\%=100-20\alpha(T)$（$\alpha$ 为归一化转化率，由 $0$ 增到 $1$），
则 $-\mathrm{d}(m\%)/\mathrm{d}T=20\,\mathrm{d}\alpha/\mathrm{d}T$——一条始终为正的
失重峰，且
\begin{equation}
\int_{T_1}^{T_2}\left[-\frac{\mathrm{d}(m\%)}{\mathrm{d}T}\right]\mathrm{d}T
=20\int_{T_1}^{T_2}\frac{\mathrm{d}\alpha}{\mathrm{d}T}\,\mathrm{d}T
=20\ \text{个百分点},
\label{eq:ch03b-dtg-area}
\end{equation}
只要 $[T_1,T_2]$ 覆盖整个失重区间，积分值与 $\beta$ 无关，且恰好等于 TG 上的
失重量。换用归一化转化率导数 $\mathrm{d}\alpha/\mathrm{d}T$ 作图，峰面积为
$100\%$ 而不是 $20$ 个百分点；两种口径都合法，但纵轴名称、峰高数值与积分
面积必须同时说明取自哪一种。图 \ref{fig:ch03b-rate} 把这个对比画成了可直接
用来训练读图习惯的形式。
\tocheck{未对照原页面图像：式 \equref{eq:ch03b-dtg-area} 是按标准
TG/DTG 关系补出的恒等式（$\alpha$ 由 $0$ 增到 $1$ 时积分必为 $1$），
源 OCR 未给出该式；图 \ref{fig:ch03b-rate} 的曲线为模型曲线，不是源书数据。}
\end{principle}

\begin{measure}
要使不同批次的数据可比，下列条件必须固定并写进实验记录（缺任何一项，
峰温都不具备可比性）：升温速率、样品量、坩埚材质与形状、载气种类与流量、
温度程序（是否含等温段）、以及空白基线的采集条件。

常用取值范围（学科通用实践，源 OCR 未列出数值，待与本实验室规程核对）：
样品量 \qty{5}{mg}--\qty{15}{mg}，做动力学时降到 \qty{2}{mg}--\qty{5}{mg}
以抑制传热滞后；升温速率 \qty{2}{K.min^{-1}}--\qty{20}{K.min^{-1}}，
求活化能时至少取四个速率；载气流量 \qty{20}{mL.min^{-1}}--\qty{50}{mL.min^{-1}}；
坩埚优先用不与样品反应的氧化铝或铂，DSC 铝坩埚只用于较低温度段。
\tocheck{未对照原页面图像：上述数值范围为学科通用量级，源 OCR 未给出。}
\end{measure}

\begin{figure}[tbp]
\centering
\begin{tikzpicture}
% 图源数据：一级速率方程 d(alpha)/dT=(A/beta)exp(-E/RT)(1-alpha)，alpha(300 K)=0，
% 四条曲线共用 E=180 kJ/mol 与同一个 A，beta=5 K/min 时 Tp=673.000 K；
% TG=100-20alpha，DTG=-dm/dT=20 d(alpha)/dT。下列坐标由
% build/r3/recompute_ch03b_dtg.py 逐点生成（证据见 build/r3/R3-C04-evidence.txt），
% 每一组都含精确峰温节点，所以画出来的峰位就是模型峰位。
\begin{groupplot}[
  group style={group size=1 by 2, vertical sep=9mm},
  xmin=523, xmax=823, xtick={550,600,650,700,750,800},
  tick label style={font={\fontsize{8}{9.6}\selectfont}},
  label style={font={\fontsize{8}{9.6}\selectfont}},
  legend style={font={\fontsize{7.5}{9}\selectfont}, draw=none, fill=none,
                cells={anchor=west}},
  grid=major, major grid style={draw=gray!22, line width=0.3pt},
  every axis plot/.append style={line width=0.9pt},
]
\nextgroupplot[width=\onecolwidth, height=5.4cm,
  ylabel={残余质量 / \%}, ymin=76, ymax=103, ytick={80,90,100},
  legend pos=south west, xticklabel pos=lower]
% TG : beta = 2.5 K/min, Tp = 659.611 K
\addplot[black] coordinates{%
(523,99.9977) (545,99.9869) (568,99.9293) (585,99.7740) (595,99.5672) (597.5,99.4929)
(600,99.4068) (602.5,99.3071) (605,99.1921) (607.5,99.0596) (610,98.9073) (612.5,98.7327)
(615,98.5331) (617.5,98.3053) (620,98.0463) (622.5,97.7526) (625,97.4209) (627.5,97.0475)
(630,96.6289) (632.5,96.1618) (635,95.6432) (637.5,95.0704) (640,94.4416) (642.5,93.7560)
(645,93.0140) (647.5,92.2176) (650,91.3705) (652.5,90.4787) (655,89.5506) (656.611,88.9381)
(657.5,88.5968) (659.611,87.7810) (660,87.6304) (662.5,86.6667) (662.611,86.6242) (665,85.7224)
(667.5,84.8152) (670,83.9624) (672.5,83.1803) (675,82.4822) (677.5,81.8777) (680,81.3716)
(682.5,80.9637) (685,80.6482) (687.5,80.4154) (690,80.2522) (692.5,80.1442) (695,80.0771)
(697.5,80.0383) (700,80.0175) (702.5,80.0073) (705,80.0028) (707.5,80.0009) (710,80.0003)
(712.5,80.0001) (715,80.0000) (717.5,80.0000) (720,80.0000) (722.5,80.0000) (725,80.0000)
(727.5,80.0000) (730,80.0000) (732.5,80.0000) (735,80.0000) (737.5,80.0000) (740,80.0000)
(742.5,80.0000) (745,80.0000) (747.5,80.0000) (750,80.0000) (752.5,80.0000) (755,80.0000)
(757.5,80.0000) (760,80.0000) (762.5,80.0000) (765,80.0000) (767.5,80.0000) (770,80.0000)
(780,80.0000) (795,80.0000) (810,80.0000) (823,80.0000)
};
% TG : beta = 5 K/min, Tp = 673.000 K
\addplot[blue!70!black,dashed] coordinates{%
(523,99.9989) (545,99.9935) (568,99.9646) (585,99.8867) (595,99.7824) (597.5,99.7448)
(600,99.7012) (602.5,99.6505) (605,99.5919) (607.5,99.5241) (610,99.4460) (612.5,99.3560)
(615,99.2526) (617.5,99.1339) (620,98.9980) (622.5,98.8428) (625,98.6660) (627.5,98.4648)
(630,98.2367) (632.5,97.9788) (635,97.6879) (637.5,97.3611) (640,96.9951) (642.5,96.5868)
(645,96.1332) (647.5,95.6317) (650,95.0801) (652.5,94.4767) (655,93.8207) (657.5,93.1124)
(660,92.3535) (662.5,91.5470) (665,90.6980) (667.5,89.8134) (670,88.9022) (672.5,87.9753)
(673,87.7892) (675,87.0458) (676,86.6764) (677.5,86.1281) (680,85.2376) (682.5,84.3901)
(685,83.6006) (687.5,82.8824) (690,82.2460) (692.5,81.6983) (695,81.2420) (697.5,80.8752)
(700,80.5919) (702.5,80.3824) (705,80.2348) (707.5,80.1364) (710,80.0744) (712.5,80.0379)
(715,80.0179) (717.5,80.0078) (720,80.0031) (722.5,80.0011) (725,80.0004) (727.5,80.0001)
(730,80.0000) (732.5,80.0000) (735,80.0000) (737.5,80.0000) (740,80.0000) (742.5,80.0000)
(745,80.0000) (747.5,80.0000) (750,80.0000) (752.5,80.0000) (755,80.0000) (757.5,80.0000)
(760,80.0000) (762.5,80.0000) (765,80.0000) (767.5,80.0000) (770,80.0000) (780,80.0000)
(795,80.0000) (810,80.0000) (823,80.0000)
};
% TG : beta = 10 K/min, Tp = 686.927 K
\addplot[red!75!black,densely dotted] coordinates{%
(523,99.9994) (545,99.9967) (568,99.9823) (585,99.9433) (595,99.8909) (597.5,99.8720)
(600,99.8500) (602.5,99.8245) (605,99.7949) (607.5,99.7606) (610,99.7211) (612.5,99.6754)
(615,99.6227) (617.5,99.5622) (620,99.4926) (622.5,99.4128) (625,99.3215) (627.5,99.2171)
(630,99.0980) (632.5,98.9625) (635,98.8085) (637.5,98.6339) (640,98.4364) (642.5,98.2136)
(645,97.9629) (647.5,97.6815) (650,97.3667) (652.5,97.0157) (655,96.6257) (657.5,96.1941)
(660,95.7185) (662.5,95.1967) (665,94.6274) (667.5,94.0096) (670,93.3433) (672.5,92.6296)
(675,91.8708) (677.5,91.0707) (680,90.2349) (682.5,89.3703) (683.927,88.8673) (685,88.4860)
(686.927,87.7976) (687.5,87.5926) (689.927,86.7280) (690,86.7022) (692.5,85.8280) (695,84.9840)
(697.5,84.1838) (700,83.4405) (702.5,82.7654) (705,82.1672) (707.5,81.6514) (710,81.2200)
(712.5,80.8708) (715,80.5985) (717.5,80.3944) (720,80.2482) (722.5,80.1484) (725,80.0839)
(727.5,80.0446) (730,80.0221) (732.5,80.0102) (735,80.0043) (737.5,80.0017) (740,80.0006)
(742.5,80.0002) (745,80.0001) (747.5,80.0000) (750,80.0000) (752.5,80.0000) (755,80.0000)
(757.5,80.0000) (760,80.0000) (762.5,80.0000) (765,80.0000) (767.5,80.0000) (770,80.0000)
(780,80.0000) (795,80.0000) (810,80.0000) (823,80.0000)
};
% TG : beta = 20 K/min, Tp = 701.424 K
\addplot[green!45!black,dash pattern=on 4pt off 2pt] coordinates{%
(523,99.9997) (545,99.9984) (568,99.9911) (585,99.9716) (595,99.9454) (597.5,99.9359)
(600,99.9249) (602.5,99.9120) (605,99.8972) (607.5,99.8800) (610,99.8600) (612.5,99.8370)
(615,99.8105) (617.5,99.7799) (620,99.7447) (622.5,99.7042) (625,99.6578) (627.5,99.6046)
(630,99.5438) (632.5,99.4743) (635,99.3951) (637.5,99.3049) (640,99.2023) (642.5,99.0859)
(645,98.9541) (647.5,98.8050) (650,98.6369) (652.5,98.4476) (655,98.2350) (657.5,97.9967)
(660,97.7305) (662.5,97.4337) (665,97.1040) (667.5,96.7389) (670,96.3360) (672.5,95.8931)
(675,95.4083) (677.5,94.8800) (680,94.3073) (682.5,93.6896) (685,93.0277) (687.5,92.3228)
(690,91.5777) (692.5,90.7963) (695,89.9840) (697.5,89.1475) (698.424,88.8338) (700,88.2952)
(701.424,87.8064) (702.5,87.4370) (704.424,86.7791) (705,86.5836) (707.5,85.7471) (710,84.9396)
(712.5,84.1733) (715,83.4597) (717.5,82.8086) (720,82.2280) (722.5,81.7229) (725,81.2955)
(727.5,80.9444) (730,80.6654) (732.5,80.4516) (735,80.2942) (737.5,80.1832) (740,80.1085)
(742.5,80.0609) (745,80.0322) (747.5,80.0159) (750,80.0074) (752.5,80.0031) (755,80.0012)
(757.5,80.0004) (760,80.0001) (762.5,80.0000) (765,80.0000) (767.5,80.0000) (770,80.0000)
(780,80.0000) (795,80.0000) (810,80.0000) (823,80.0000)
};
\legend{\qty{2.5}{K.min^{-1}}, \qty{5}{K.min^{-1}},
        \qty{10}{K.min^{-1}}, \qty{20}{K.min^{-1}}}
\nextgroupplot[width=\onecolwidth, height=5.4cm,
  xlabel={温度 / K},
  ylabel={$-\mathrm{d}m/\mathrm{d}T$ / \%\,K$^{-1}$（正峰表示失重）},
  ymin=0, ymax=0.51, ytick={0,0.1,0.2,0.3,0.4,0.5}, legend pos=north west]
% DTG : beta = 2.5 K/min, Tp = 659.611 K
\addplot[black] coordinates{%
(523,0.00019) (545,0.00100) (568,0.00498) (585,0.01496) (595,0.02758) (597.5,0.03199)
(600,0.03704) (602.5,0.04280) (605,0.04936) (607.5,0.05679) (610,0.06520) (612.5,0.07466)
(615,0.08528) (617.5,0.09714) (620,0.11030) (622.5,0.12484) (625,0.14079) (627.5,0.15816)
(630,0.17692) (632.5,0.19696) (635,0.21814) (637.5,0.24022) (640,0.26285) (642.5,0.28559)
(645,0.30788) (647.5,0.32903) (650,0.34824) (652.5,0.36461) (655,0.37718) (656.611,0.38281)
(657.5,0.38498) (659.611,0.38717) (660,0.38709) (662.5,0.38277) (662.611,0.38241) (665,0.37149)
(667.5,0.35313) (670,0.32798) (672.5,0.29684) (675,0.26101) (677.5,0.22225) (680,0.18260)
(682.5,0.14415) (685,0.10887) (687.5,0.07826) (690,0.05326) (692.5,0.03411) (695,0.02041)
(697.5,0.01133) (700,0.00579) (702.5,0.00270) (705,0.00114) (707.5,0.00043) (710,0.00014)
(712.5,0.00004) (715,0.00001) (717.5,0.00000) (720,0.00000) (722.5,0.00000) (725,0.00000)
(727.5,0.00000) (730,0.00000) (732.5,0.00000) (735,0.00000) (737.5,0.00000) (740,0.00000)
(742.5,0.00000) (745,0.00000) (747.5,0.00000) (750,0.00000) (752.5,0.00000) (755,0.00000)
(757.5,0.00000) (760,0.00000) (762.5,0.00000) (765,0.00000) (767.5,0.00000) (770,0.00000)
(780,0.00000) (795,0.00000) (810,0.00000) (823,0.00000)
};
% DTG : beta = 5 K/min, Tp = 673.000 K
\addplot[blue!70!black,dashed] coordinates{%
(523,0.00009) (545,0.00050) (568,0.00250) (585,0.00752) (595,0.01394) (597.5,0.01620)
(600,0.01880) (602.5,0.02178) (605,0.02519) (607.5,0.02909) (610,0.03353) (612.5,0.03857)
(615,0.04430) (617.5,0.05077) (620,0.05806) (622.5,0.06625) (625,0.07543) (627.5,0.08566)
(630,0.09701) (632.5,0.10955) (635,0.12333) (637.5,0.13837) (640,0.15466) (642.5,0.17218)
(645,0.19084) (647.5,0.21049) (650,0.23093) (652.5,0.25186) (655,0.27291) (657.5,0.29360)
(660,0.31335) (662.5,0.33148) (665,0.34725) (667.5,0.35985) (670,0.36843) (672.5,0.37219)
(673,0.37231) (675,0.37045) (676,0.36809) (677.5,0.36268) (680,0.34862) (682.5,0.32836)
(685,0.30236) (687.5,0.27152) (690,0.23714) (692.5,0.20082) (695,0.16434) (697.5,0.12949)
(700,0.09783) (702.5,0.07056) (705,0.04834) (707.5,0.03129) (710,0.01902) (712.5,0.01078)
(715,0.00566) (717.5,0.00273) (720,0.00120) (722.5,0.00048) (725,0.00017) (727.5,0.00005)
(730,0.00001) (732.5,0.00000) (735,0.00000) (737.5,0.00000) (740,0.00000) (742.5,0.00000)
(745,0.00000) (747.5,0.00000) (750,0.00000) (752.5,0.00000) (755,0.00000) (757.5,0.00000)
(760,0.00000) (762.5,0.00000) (765,0.00000) (767.5,0.00000) (770,0.00000) (780,0.00000)
(795,0.00000) (810,0.00000) (823,0.00000)
};
% DTG : beta = 10 K/min, Tp = 686.927 K
\addplot[red!75!black,densely dotted] coordinates{%
(523,0.00005) (545,0.00025) (568,0.00125) (585,0.00377) (595,0.00701) (597.5,0.00815)
(600,0.00947) (602.5,0.01099) (605,0.01273) (607.5,0.01472) (610,0.01700) (612.5,0.01961)
(615,0.02257) (617.5,0.02595) (620,0.02979) (622.5,0.03413) (625,0.03904) (627.5,0.04457)
(630,0.05080) (632.5,0.05777) (635,0.06557) (637.5,0.07425) (640,0.08389) (642.5,0.09453)
(645,0.10624) (647.5,0.11905) (650,0.13297) (652.5,0.14802) (655,0.16415) (657.5,0.18130)
(660,0.19935) (662.5,0.21813) (665,0.23740) (667.5,0.25686) (670,0.27611) (672.5,0.29470)
(675,0.31207) (677.5,0.32760) (680,0.34062) (682.5,0.35042) (683.927,0.35430) (685,0.35630)
(686.927,0.35775) (687.5,0.35762) (689.927,0.35402) (690,0.35383) (692.5,0.34458) (695,0.32975)
(697.5,0.30950) (700,0.28436) (702.5,0.25515) (705,0.22305) (707.5,0.18944) (710,0.15587)
(712.5,0.12382) (715,0.09463) (717.5,0.06931) (720,0.04843) (722.5,0.03214) (725,0.02015)
(727.5,0.01186) (730,0.00652) (732.5,0.00332) (735,0.00156) (737.5,0.00067) (740,0.00026)
(742.5,0.00009) (745,0.00003) (747.5,0.00001) (750,0.00000) (752.5,0.00000) (755,0.00000)
(757.5,0.00000) (760,0.00000) (762.5,0.00000) (765,0.00000) (767.5,0.00000) (770,0.00000)
(780,0.00000) (795,0.00000) (810,0.00000) (823,0.00000)
};
% DTG : beta = 20 K/min, Tp = 701.424 K
\addplot[green!45!black,dash pattern=on 4pt off 2pt] coordinates{%
(523,0.00002) (545,0.00013) (568,0.00062) (585,0.00189) (595,0.00351) (597.5,0.00409)
(600,0.00475) (602.5,0.00552) (605,0.00640) (607.5,0.00740) (610,0.00856) (612.5,0.00988)
(615,0.01139) (617.5,0.01312) (620,0.01509) (622.5,0.01732) (625,0.01986) (627.5,0.02274)
(630,0.02599) (632.5,0.02967) (635,0.03381) (637.5,0.03846) (640,0.04369) (642.5,0.04953)
(645,0.05605) (647.5,0.06331) (650,0.07135) (652.5,0.08024) (655,0.09002) (657.5,0.10074)
(660,0.11243) (662.5,0.12512) (665,0.13880) (667.5,0.15345) (670,0.16902) (672.5,0.18543)
(675,0.20253) (677.5,0.22016) (680,0.23808) (682.5,0.25598) (685,0.27350) (687.5,0.29021)
(690,0.30561) (692.5,0.31916) (695,0.33028) (697.5,0.33835) (698.424,0.34044) (700,0.34279)
(701.424,0.34350) (702.5,0.34308) (704.424,0.34020) (705,0.33879) (707.5,0.32964) (710,0.31555)
(712.5,0.29670) (715,0.27354) (717.5,0.24677) (720,0.21738) (722.5,0.18653) (725,0.15552)
(727.5,0.12563) (730,0.09801) (732.5,0.07361) (735,0.05302) (737.5,0.03648) (740,0.02387)
(742.5,0.01478) (745,0.00862) (747.5,0.00470) (750,0.00239) (752.5,0.00112) (755,0.00048)
(757.5,0.00019) (760,0.00007) (762.5,0.00002) (765,0.00001) (767.5,0.00000) (770,0.00000)
(780,0.00000) (795,0.00000) (810,0.00000) (823,0.00000)
};
\legend{\qty{2.5}{K.min^{-1}}, \qty{5}{K.min^{-1}},
        \qty{10}{K.min^{-1}}, \qty{20}{K.min^{-1}}}
\end{groupplot}
\end{tikzpicture}
\caption{不同升温速率下 TG 与 DTG 的模型曲线（\modelmark，非原书数据）。
两块面板共用横轴（单位为 K，与式~\eqref{eq:ch03b-kissinger}
式~\eqref{eq:ch03b-dTp} 的 $T_{\mathrm p}\approx\qty{673}{K}$ 同一单位）。
读者应看三件事：(a) 四条 TG 曲线的两个平台都落在 100\% 与 80\%，只有转变
区间随升温速率右移；(b) DTG 峰顶由 659.6 K 移到 701.4 K（即 386.4 \degC{} 到
428.3 \degC），升温速率变化八倍峰温移动 41.8 K，每翻一倍上移 13.4--14.5 K；
(c) 纵轴取实际质量导数 $-\mathrm{d}m/\mathrm{d}T$（正峰表示失重）后，四条 DTG
曲线的峰面积都是 $20$ 个质量百分点，等于 TG 上的失重量，与升温速率无关；
但峰高并不严格相等（依次为 $0.387$、$0.372$、$0.358$、$0.343$
\%\,K$^{-1}$，随升温速率略降），因为峰温升高使 $E/(R\,T_{\mathrm p}^{2})$ 变小。
取归一化转化率导数 $\mathrm{d}\alpha/\mathrm{d}T$ 作图时，峰高改为上列数值的
$5$ 倍（$1.94$、$1.86$、$1.79$、$1.72$ \%\,K$^{-1}$），面积改为 $100\%$，
结论不变；两种口径不能混用。
模型条件：四条曲线由同一条一级速率方程
$\mathrm{d}\alpha/\mathrm{d}T=(A/\beta)\exp[-E/(R\,T)]\,(1-\alpha)$ 积分得到，
固定失重 20\%，共用 $E=\qty{180}{kJ.mol^{-1}}$ 与同一个频率因子 $A$，
峰温由动力学自动给出（$\beta=\qty{5}{K.min^{-1}}$ 时 $T_{\mathrm p}=673.0$ K，
其余三条由同一组参数解出）；按式~\eqref{eq:ch03b-kissinger} 回算这四组峰温
得到的 $E$ 与设定值一致。峰温处的转化率 $\alpha(T_{\mathrm p})\approx0.61$，
所以峰温不能当作半转化温度。温度程序自 300 K 起积分，523 K 处
$\alpha<2\times10^{-4}$，因此两条平台在图上仍读作 100\% 与 80\%。
真实样品的 DTG 峰高一般随升温速率上升（实际失重速率还随 $\beta$ 增大）；
这里用同一坐标作图是为了让「峰面积不变」直接可见。未计入传热滞后、浮力与
基线漂移。
\figsrc{据本章动力学讨论绘制，上册 §3.3，书内 p.181 / PDF p.200。}}
\label{fig:ch03b-rate}
\end{figure}

\begin{example}[读图：升温速率改变的是峰温还是峰面积]
\begin{obs}
在图 \ref{fig:ch03b-rate} 上逐条读出的内容：
(a) 面板中，升温速率从 \qty{2.5}{K.min^{-1}} 增到 \qty{20}{K.min^{-1}} 时，
四条曲线的起始平台都在 100\%，终止平台都在 80\%，转变段的中心依次右移；
(b) 面板中，DTG 峰顶的位置依次为 659.6、673.0、686.9、701.4 K（即 386.4、
399.9、413.8、428.3 \degC），八倍速率之间峰温移动 41.8 K；把这四组峰温按
式~\equref{eq:ch03b-kissinger} 作
$\ln(\beta/T_{\mathrm p}^{2})$--$1/T_{\mathrm p}$ 图，回算得到的 $E$ 就是
生成曲线时设定的 \qty{180}{kJ.mol^{-1}}；
(c) 峰的最大高度（纵轴为实际质量导数 $-\mathrm{d}m/\mathrm{d}T$）依次为
0.387、0.372、0.358、0.343 \%\,K$^{-1}$，随升温速率略降；四个峰的半高宽
依次约为 47、49、51、53 K（约为转变区间宽度的三分之一）。峰温处 TG 的
残余质量为 87.8\%（$\beta=\qty{5}{K.min^{-1}}$），即该处 $\alpha\approx0.61$，
不是一半。
两步之间没有出现新的失重台阶，也没有出现肩峰。
\end{obs}
\begin{infer}
由「两个平台位置相同」推断：这一步的总失重量为 20\%，与升温速率无关，
因此失重量可以作为定量的依据。由「峰顶随升温速率单调右移、且位移量随
$\ln(\beta_2/\beta_1)$ 近似线性」推断：峰温位移来自动力学，不是材料在
不同实验里变成了不同的东西。由「只有一个峰、无肩峰」推断：至少在当前的
分辨能力下，这一步可以按单一过程处理。

不能由本图推断的内容：不能推断该过程的活化能具体是多少——图
\ref{fig:ch03b-rate} 的 $E=\qty{180}{kJ.mol^{-1}}$ 是生成曲线时\emph{设定}的
参数，上一条回算只是自洽性检验，不是对真实样品的测量；真实体系必须用实测
峰温按式~\equref{eq:ch03b-kissinger} 拟合才能得到 $E$；
不能推断样品的热稳定性温度（峰温不是稳定性判据）；也不能排除存在两个
活化能相近、峰位重叠的步骤。
\end{infer}
\end{example}

\begin{dataflow}
从原始曲线到可比的量，四步：

① 扣空白与浮力校正。同一温度程序、同一坩埚、同一流量下先测空坩埚，
从样品曲线中减去；未做这一步，峰位判读本身就不成立。

② 定位。用 DTG 峰顶定 $T_{\mathrm p}$，用平台差定失重量 $\Delta m$，
用 DTG 的起始偏离点与回折点定温度区间。

③ 判别位移来源。把 $T_{\mathrm p}$ 按式 \equref{eq:ch03b-kissinger} 作
$\ln(\beta/T_{\mathrm p}^{2})$ 对 $1/T_{\mathrm p}$ 的图：若成直线，
位移全部来自动力学；若明显偏离直线，说明发生了传热控制或反应路径改变。

④ 定量。用失重量而不是峰高做定量；只在步骤 ①--③ 都通过后才比较不同样品。
\end{dataflow}

\begin{table}[tbp]
\centering
\caption{实验条件对曲线的三种作用：峰温、峰形、峰面积。读者应看出第 3、4 列
的规律——峰温几乎对所有条件敏感，峰面积只对「反应路径改变」与「基线未扣」
敏感。表中「基本不变」指该量为守恒量时的不变性，仍可能被重叠峰拆分误差放大。
\src{据上册 §3.6，书内 p.209—212 / PDF p.228—231 归纳；表中峰温位移量按
Arrhenius 动力学估算，非源书实验数据。}}
\label{tab:ch03b-conditions}
\small
\begin{tabular}{@{}p{42mm}p{28mm}p{28mm}p{38mm}@{}}
\toprule
实验因素（作用机制） & 对峰温 $T_{\mathrm p}$ & 对峰形 & 对峰面积／失重量 \\
\midrule
升温速率 $\beta$ 增大（动力学控制，式 \equref{eq:ch03b-kissinger}） &
显著上移，八倍速率可移 40 K 以上 &
整体向高温展宽，相邻峰更易重叠 &
基本不变 \\
样品量增大（颗粒床内部传热与传质滞后，自生气氛增强） &
上移 &
展宽、拖尾，分辨下降 &
名义不变，但重叠峰拆分误差被放大 \\
气氛改变（惰性转为氧化性） &
可大幅移动，甚至原峰消失、新峰出现 &
峰数改变（出现燃烧峰） &
\textbf{会真正改变}：反应路径不同，失重量不同 \\
载气流量增大（产物分压降低，抑制逆反应） &
略降 &
略变窄 &
一般不变；流量突变会引起假失重或假增重 \\
坩埚材质与形状改变（催化、与样品反应、热导与热容不同） &
改变 &
可能畸变，可出现假峰 &
坩埚参与反应时失真 \\
粒度与堆积改变（大颗粒内部温度梯度，堆积层厚） &
大颗粒、厚堆积时上移 &
展宽 &
基本不变 \\
浮力与基线未扣（气体密度随温度变化） &
表观漂移，峰位判读失真 &
整体倾斜，出现假台阶 &
直接污染，产生假失重或假增重 \\
\bottomrule
\end{tabular}

\vspace{1.5mm}
{\footnotesize 结论：\textbf{峰面积（总失重量、总焓变）是条件相对稳健的量，
峰温是最容易被条件左右的量}。任何「用峰温比较材料稳定性」的结论，都必须
先固定升温速率、样品量、气氛、坩埚与流量这五项。}
\end{table}

\begin{pitfall}
\textbf{把峰温当材料常数。} 由式 \equref{eq:ch03b-dTp}，\qty{5}{K.min^{-1}}
下测得的 \qty{420}{\celsius} 与 \qty{20}{K.min^{-1}} 下测得的
\qty{455}{\celsius} 完全可以是同一种材料。跨文献比较峰温之前，必须先确认
升温速率相同。

\textbf{用 DTG 峰高代替峰面积。} 峰高正比于升温速率（$\mathrm{d}m/\mathrm{d}t
= \beta\,\mathrm{d}m/\mathrm{d}T$），换一个速率峰高就变，峰面积才守恒。
图 \ref{fig:ch03b-rate} 的纵轴取 $-\mathrm{d}m/\mathrm{d}T$ 而不是
$-\mathrm{d}m/\mathrm{d}t$，就是为了让四个速率下的曲线落在同一坐标上比较；
换成 $-\mathrm{d}m/\mathrm{d}t$ 作图，峰高会随 $\beta$ 按比例升高。

\textbf{把肩峰当新相。} 样品量偏大或粒度不均时，同一个过程会因床层温度
分布而裂成主峰加肩峰。判据：把样品量减半，若肩峰消失或大幅减弱，
则它来自传热而非新相。
\end{pitfall}

\begin{crosscheck}
三项回查按顺序做：① 空白基线回查——空坩埚曲线在同一温度区间应当平坦，
若不平坦，先修基线；② 样品量回查——样品量减半后重测，峰温应下移或不变，
若显著上移则原曲线受传热控制；③ 升温速率回查——用两个以上速率重测，
按式 \equref{eq:ch03b-kissinger} 作图，直线性本身就是数据自洽的证据。
三项都通过后，才允许把峰温写进结论，并且必须与升温速率一起写。
\end{crosscheck}

\subsection{气氛：改变的是反应路径，不只是反应速率}

\begin{probout}
气氛不是「背景条件」，而是反应物或产物的一部分。同一样品在空气与氮气下
得到的两条曲线，其失重步数、失重温度区间与失重量都可能不同，因此
\textbf{不可互相比较，也不能互相换算}。本节要回答：如何从曲线形状判断
气氛是否改变了反应路径，以及如何在两种气氛之间建立可比的实验设计。
\end{probout}

\begin{principle}
惰性气氛下发生的是热解、脱羟基、脱碳酸盐一类反应，固相自身分解，
逸出物是分解产物；氧化性气氛下，同样的有机或还原性组分会被氧化，
逸出物是 \ce{CO2} 与 \ce{H2O}，同时释放热量。

这一差别在两类曲线上的表现是互补的：
\begin{itemize}
  \item TG 上，惰性气氛下碳质残留会形成残炭，失重不彻底；空气下残炭被
  氧化烧掉，失重量更大，且多出一个高温失重步。因此\textbf{用空气下的
  失重量测碳含量是可行的，用氮气下的失重量测碳含量则系统性偏低}。
  \item DSC/DTA 上，同一过程在惰性下常表现为吸热（分解需要能量），
  在氧化性下可能表现为放热（燃烧放热）。峰的热效应符号可以直接翻转，
  这是判断「这一步是不是燃烧」最直接的证据。
\end{itemize}
本书中积炭催化剂的 DSC 测试正是用这一判别方法：先在惰性气氛下升温，
切换到氧气后出现一个明显的放热峰，该峰被归因于催化剂表面积炭被氧化成
\ce{CO2}、\ce{CO} 等挥发物；由此得到「积炭覆盖活性中心或堵塞孔道导致
活性下降」的结论。\src{上册 §3.4.1，书内 p.184--185 / PDF p.203--204}
这里读出的是两件事：失重步的出现是质量证据，放热峰的出现是热效应证据，
两者合起来才支持「可燃的碳质物种」这一归属；仅凭其中一个都会过度解读。
\end{principle}

\begin{measure}
气氛条件的五项记录：气体种类与纯度、流量、样品处实际气氛（是否被自身
分解产物改性）、是否中途切换气氛、切换时的温度与稳定时间。

三个容易被忽略的实际问题：
\begin{itemize}
  \item \textbf{残余氧。} 高纯氮中若含氧量偏高，碳的氧化峰依然会出现，
  但位置和形状与空气下不同，容易误判为「材料本身的低温氧化」。
  怀疑时用氩气复测或检测尾气含氧量。
  \item \textbf{自生气氛。} 样品量偏大时，分解产物在床层内累积，
  使样品实际处于还原性或产物富集的气氛中，与名义气氛不同。
  这是样品量与气氛两个变量的耦合，减小样品量即可验证。
  \item \textbf{气体的热导率影响 DSC 基线。} 切换气氛会改变基线，
  所以两种气氛下的曲线必须各自做各自的空白基线，不能共用。
\end{itemize}
\tocheck{未对照原页面图像：本段的残余氧与自生气氛讨论为学科通用实践，
源 OCR 未展开；切换气氛的实例见书内 p.184--185 的 DSC 测试描述。}
\end{measure}

\begin{dataflow}
判断「气氛是否改变了反应路径」的四步：

① 在惰性气氛与空气下各做一次，其余条件（速率、样品量、坩埚、流量）完全相同；

② 比较失重步数：空气下多出的失重步通常是可燃组分的氧化；

③ 比较热效应符号：惰性下吸热、氧化下放热的同温区过程，基本可以判定为
氧化反应；

④ 比较总失重量：若两者差值与残炭量相当，则「惰性下残炭、空气下烧尽」
的解释成立；若差值与残炭量不符，说明还有别的路径，需要逸出物分析。
\end{dataflow}

\begin{pitfall}
\textbf{跨气氛比较特征温度。} 把文献中氮气下的分解温度与本实验室空气下的
起始氧化温度并列，得到的「稳定性对比」没有意义：空气下测的是氧化起始温度，
氮气下测的是热分解温度，两者不属于同一个过程。

\textbf{把气氛效应当催化效应。} 掺杂样品在空气下氧化温度降低，可能只是
掺杂提高了氧的传输或增加了比表面积，而不是催化了分解本身。区分方法是
在惰性气氛下重测同一对样品：若惰性下的峰温不降，则空气下的降低属于氧化
路径的促进。

\textbf{忘记气氛流量对峰温的次级影响。} 流量增大使产物分压降低，
对可逆分解过程会略微降低峰温。在同一气氛内比较不同样品时影响不大，
但把「流量不同的两组数据」当作「气氛不同的两组数据」来解读则会出错。
\end{pitfall}

\begin{crosscheck}
气氛结论至少需要一项非热分析证据：元素分析（CHN）给出碳含量、
XRD 给出固相物相、尾气分析给出逸出物组成。对积炭一类问题，
「空气下的失重增量」应当与元素分析测得的碳含量在误差范围内一致。
本书中 \ce{CuO}-\ce{ZnO}-\ce{Al2O3}-\ce{CrO_x} 前驱体组成的确定，
用的就是这种质量平衡回查（见下节实例），只是那里的对照对象是
按化学计量计算的理论失重。
\end{crosscheck}

\subsection{五类典型误判与回查方法}

\begin{probout}
本节把热分析中最常见、也最容易写进论文结论的五类误判集中列出，
每一条都给出「误判长什么样」与「用什么实验回查」。这一节的用法不是通读，
而是在得到一条反常结论时逐条对照。
\end{probout}

\begin{principle}
五类误判有一个共同的根源：把\textbf{峰温}当成了材料的固有属性。
峰温由动力学与传热共同决定（式 \equref{eq:ch03b-kissinger}、
式 \equref{eq:ch03b-dTp}），只有峰面积与失重量属于守恒量。凡是从峰位
出发的结论，都需要额外证据；凡是从峰面积出发的结论，只需要确认基线正确
与过程单一。
\end{principle}

\begin{pitfall}
\textbf{（一）把峰温位移当材料稳定性差异。} 两组样品的 DTG 峰温相差 30 K，
结论写成「B 比 A 稳定」。真实原因常常只是两次实验的升温速率不同
（\qty{5}{K.min^{-1}} 对 \qty{10}{K.min^{-1}}），或样品量相差一倍。

\emph{回查}：固定升温速率重测；或对每个样品各取四个升温速率，
按式 \equref{eq:ch03b-kissinger} 求 $E$，用 $E$ 而不是 $T_{\mathrm p}$
比较稳定性。本书中 12-钨硅酸负载到 HEMT 沸石后分解放热峰由
\qty{522.12}{\celsius} 移到 \qty{559.13}{\celsius}，
结论是「二者之间发生了相互作用」，这一比较之所以成立，
是因为两条件在同一升温速率下测得。\src{上册 §3.4.1，
书内 p.183--184 / PDF p.202--203} 反过来，如果这两个数字来自不同的
升温速率，则该结论不成立。

\textbf{（二）基线未扣或浮力未校正。} 表现为低温段缓慢漂移、高温段整体
倾斜，或出现一个「谁也没见过」的宽失重台阶。浮力随温度升高使表观质量
单调变化，其大小与升温速率和气体种类都有关，常被误读为「表面吸附水的
连续脱除」。

\emph{回查}：同一条件下测空坩埚，用同一程序做空白并相减；
再把载气由氮气换成氩气，真正的样品失重不变，浮力引起的漂移会变。

\textbf{（三）坩埚与样品反应。} 坩埚不是惰性容器。碱金属或碱土金属的
碳酸盐、硫酸盐在高温下可与氧化铝或石英坩埚反应；铂坩埚对有机物的
氧化有催化作用，也可能与某些金属形成合金；镍、铁一类样品在高温下
与氧化铝坩埚发生界面反应。表现是峰温偏移、出现附加峰，
严重时失重量也失真。

\emph{回查}：换一种材质的坩埚（如氧化铝换铂或换石墨）重测同一份样品。
若曲线不变，坩埚惰性成立；若曲线改变，先前所有定量结论都需要重做。
\tocheck{未对照原页面图像：坩埚与样品反应的具体体系为学科通用认识，
源 OCR 未列出实例，待与本实验室规程核对。}

\textbf{（四）样品量过大致传热滞后。} 床层内部存在温度梯度，
热电偶测的是坩埚附近的温度而不是样品中心的温度，于是样品实际温度低于
记录温度，曲线向高温方向整体平移并展宽。表现是峰温偏高、相邻峰合并、
肩峰出现。对动力学分析而言，这会把表观活化能算小。

\emph{回查}：样品量减半（必要时再减半）重测。若峰温下移且峰变窄，
说明原曲线受传热控制，应取小样品量的数据；若峰温不变，样品量不是变量。

\textbf{（五）用 TG 失重直接推断物种而不做逸出物分析。}
失重 18\% 就被写成「失去两个结晶水」，失重 6\% 就被写成「脱除结构羟基」。
数值巧合不能作为物证：不同的物种组合可以给出同一个失重量，
同一步失重也可能同时包含多个物种。

\emph{回查}：三件事至少做一件——按化学计量计算\textbf{理论失重}并与实测比较；
用联用技术做逸出物分析；用 XRD 或元素分析确认终态固相。
本书中前驱体组成的确定正是这一思路的正例：先按可能的沉淀物组成算出理论
分解失重，与实测值比较，发现理论值系统性大于实测值且偏差较大，
据此判断沉淀物的设计有问题；改用铬酸钠作铬源重新设计沉淀反应后，
理论失重与实测失重的偏差缩小到 \qty{0.9}{\percent} 以内，
才确认新组成正确。\src{上册 §3.4.1，书内 p.186--188 / PDF p.205--207}
这里 TG 的作用是\emph{验证}组成假设，不是\emph{发现}组成。

另一个反例是球磨 \ce{Al}-\ce{SnCl2} 复合物的 DSC 测试：第一次升温出现两个
吸热峰，第二次升温第一个峰消失、第二个峰保留，说明前者不可逆、后者可逆；
可逆峰的起始点 \qty{225.2}{\celsius} 与纯锡熔点 \qty{230}{\celsius} 接近，
不可逆峰的起始点 \qty{180.9}{\celsius} 与 \ce{AlCl3} 升华温度
\qty{183}{\celsius} 接近。若只看 DSC 曲线本身，这两个峰最多只能说成
「某物熔化／升华」；只有把 XRD 检出新相金属锡、EDS 检出氯元素分布，
才能写成「\ce{Sn^2+} 被还原为金属锡并生成 \ce{AlCl3}」。
\src{上册 §3.4.4，书内 p.195--196 / PDF p.214--215}
\end{pitfall}

\begin{crosscheck}
把五条回查整理成一张可执行的清单：

\begin{itemize}
  \item 峰位类结论 → 用「固定升温速率重测」+「Kissinger 图直线性」回查；
  \item 失重量类结论 → 用「空白基线」+「理论化学计量失重」回查；
  \item 物种归属类结论 → 用「逸出物分析」+「XRD／元素分析」回查；
  \item 稳定性类结论 → 用「多个升温速率」+「同一样品量」回查；
  \item 催化效应类结论 → 用「惰性与氧化两种气氛」回查。
\end{itemize}

五条中最容易被跳过的是第一条。建议在实验记录模板里把升温速率与样品量
写成与样品名称并列的必填字段，从流程上消除这个误判。
\end{crosscheck}

\subsection{应用实例选读：从曲线到结论}

\begin{probout}
本节按「观测—推断—回查」三段式读四个实例，覆盖本书第 3 章应用部分的主要
类型：焙烧温度的确定、负载后峰温位移的解读、前驱体组成的质量平衡验证、
以及由峰温求脱氢活化能。读这些实例的重点不是记住结论，
而是看清每条结论背后必须有的独立证据。
\end{probout}

\subsubsection{焙烧温度的确定：热分析与 XRD 的互补}

\begin{dataflow}
\textbf{观测。} 含铜、锌、镍硝酸盐的沉淀混合物经干燥后进行 DTA-TG 测试
（\src{上册 §3.4.1，书内 p.182 / PDF p.201}）：\qty{180}{\celsius} 以前有
少量失重；\qty{180}{\celsius}--\qty{420}{\celsius} 之间 DTA 出现明显峰，
同时 TG 上有较大失重；温度达到 \qty{400}{\celsius} 以上，TG 与 DTA 曲线
都接近直线。

\textbf{推断。} 低温段少量失重归因于表面水分的脱除；中温段的
「DTA 峰 + TG 失重」同时出现在同一温区，说明该区间发生了沉淀物的分解；
\qty{400}{\celsius} 以上两条曲线都变平，说明分解已经完成，
故焙烧温度应取 \qty{400}{\celsius} 左右。

\textbf{回查。} 在 \qty{320}{\celsius}、\qty{400}{\celsius}、
\qty{500}{\celsius} 三个温度下焙烧后做 XRD：\qty{320}{\celsius} 时物相
基本为非晶型，衍射峰不明显；\qty{400}{\celsius} 时几种氧化物的衍射峰较为
明显；\qty{500}{\celsius} 时衍射结果与 \qty{400}{\celsius} 几乎相同，
但温度过高易出现烧结，影响比表面积与强度。综合两者，最佳焙烧温度定在
\qty{400}{\celsius} 左右。\src{上册 §3.4.1，书内 p.183 / PDF p.202}

这里的分工值得注意：热分析指出「分解在哪个温区完成」，
XRD 指出「分解后是否形成了需要的物相」。二者缺一，
「\qty{400}{\celsius} 最佳」这个结论都不完整——单看 TG 只能说分解完成，
单看 XRD 只能说物相出现，都没有给出温度下限与上限的组合判据。
\end{dataflow}

\subsubsection{负载后峰温向高温移动：相互作用还是实验条件}

\begin{dataflow}
\textbf{观测。} HEMT 沸石样品仅在 \qty{100}{\celsius} 附近出现一个脱水峰；
12-钨硅酸本身有三个峰，\qty{61.3}{\celsius} 与 \qty{191.4}{\celsius} 两个
吸热峰分别对应失去结晶水与结构水，\qty{522.12}{\celsius} 的放热峰对应
杂多酸结构的分解。HEMT 负载约三分之一杂多酸后，在 \qty{559.13}{\celsius}
出现一个小的放热峰。\src{上册 §3.4.1，书内 p.183--184 / PDF p.202--203}

\textbf{推断。} 该放热峰与 12-钨硅酸的分解峰对应，但峰位向高温方向移动，
说明负载到 HEMT 表面后杂多酸与载体之间发生了相互作用，
使其分解需要更高的温度。

\textbf{边界。} 这一推断成立的前提是两次测试的升温速率相同。
若升温速率不同，按式 \equref{eq:ch03b-dTp}，八倍速率差即可造成 40 K 以上的
峰温位移，恰好覆盖本例中 37 K 的位移量，结论就会完全反转。
因此这类比较必须在实验部分明确写出升温速率，并在图中标注。源 OCR 未给出
该实验的升温速率，此处不能替它补齐。\tocheck{未对照原页面图像：
图 3-12 的升温速率在 OCR 文本中缺失，本例的峰温可比性依赖这一条件。}
\end{dataflow}

\subsubsection{前驱体组成：用理论失重做质量平衡验证}

\begin{dataflow}
\textbf{观测。} 按投料比设计的 \ce{CuO}-\ce{ZnO}-\ce{Al2O3}-\ce{CrO_x}
催化剂前驱体，在空气下测得的 DTG 曲线呈「双馒头峰」，分解温区约为
\qty{140}{\celsius}--\qty{360}{\celsius}，到 \qty{450}{\celsius} 结束；
按沉淀物组成计算的理论热分解失重与实测失重相比，理论值均大于实测值，
偏差较大。\src{上册 §3.4.1，书内 p.185--186 / PDF p.204--205}

\textbf{推断。} 理论失重与实测失重的系统性偏差，说明所假定的沉淀物组成
不对，即前驱体的沉淀物设计有问题。

\textbf{回查。} 分别制备以 \ce{Na2CO3} 为沉淀剂和以 \ce{Na2CrO4}（铬酸钠）
为铬源的 \ce{CuO}-\ce{Al2O3} 前驱体并比较其 DTG：前者是一个尖锐的峰，
主要来自碳酸铜类沉淀的分解；后者出现的是另一个特征的峰。
据此判断，\ce{CuO}-\ce{ZnO}-\ce{Al2O3}-\ce{CrO_x} 前驱体的双馒头峰与
铬酸钠参与沉淀有关，即铬酸钠使碳酸铜不能生成或只生成氢氧化铜。
按新组成重新计算理论失重后，与实测值的偏差缩小到 \qty{0.9}{\percent}
以内（\src{上册 §3.4.1，书内 p.187--188 / PDF p.206--207}），
从而确认后一种沉淀物种的设计是正确的，并进一步说明：
以铬酸钠作铬源引入时会与硝酸铜生成铬酸铜沉淀，同时抑制碳酸铜的生成。
\end{dataflow}

\begin{pitfall}
本例是整个应用部分最容易被误读的一处。TG 曲线本身只给出「失了多少」，
它无法直接说出「失去的是什么」。这里之所以能判断组成，
是因为用了\textbf{两个独立的约束}：一是理论化学计量失重，
二是不同沉淀剂条件下 DTG 峰形的差异。如果把「双馒头峰」直接读成
「两种碳酸盐分解」，就跳过了全部验证步骤，属于上文第（五）类误判。
\end{pitfall}

\subsubsection{由峰温求脱氢活化能：峰温位移的正确用法}

\begin{dataflow}
\textbf{观测。} 对球磨处理及掺杂不同尺度 \ce{TiH2} 的 \ce{LiAlH4} 样品做
非等温脱氢的 TGA 与 DSC 测试，得到两步脱氢过程。由不同升温速率下的 DSC
峰温，按动力学方法算得两步脱氢的表观活化能：球磨样品为
\qty{96.0}{kJ.mol^{-1}} 与 \qty{92.7}{kJ.mol^{-1}}；
掺入微米级 \ce{TiH2} 后降为 \qty{82.9}{kJ.mol^{-1}} 与
\qty{82.5}{kJ.mol^{-1}}；掺入纳米级 \ce{TiH2} 后进一步降为
\qty{76.1}{kJ.mol^{-1}} 与 \qty{68.5}{kJ.mol^{-1}}。
\src{上册 §3.4.4，书内 p.196--197 / PDF p.215--216}
\tocheck{未对照原页面图像：上列活化能数值取自表 3-11 的 OCR 文本，
其中个别数字的末位存在识别风险，数值本身仅供量级判断。}

\textbf{推断。} 掺杂使两步脱氢的表观活化能同时降低，且纳米级 \ce{TiH2}
的降低幅度更大，因此纳米级 \ce{TiH2} 对 \ce{LiAlH4} 脱氢的催化效果更好；
这与 TGA 上观测到的起始脱氢温度下降一致。

\textbf{为什么这里可以用峰温。} 与上一节「不能直接用峰温比较稳定性」并不
矛盾：这里不是在比较两个温度数字，而是把\textbf{同一样品在多个升温速率下}
的峰温整体作为输入，用动力学关系提取出一个条件无关的量——活化能。
峰温在这里是中间量，不是结论。这正是误判第（一）条给出的正确做法。
\end{dataflow}

\begin{crosscheck}
四个实例的共同结构：热分析给出温度区间与质量／热效应变化，
另外一项独立测量（XRD、元素分析、理论化学计量、多速率动力学分析）
给出归属或守恒量，两者都成立才写结论。

对本节内容做回查时按下列顺序：先确认升温速率与样品量已记录，
再确认基线已扣，然后确认失重步与理论化学计量或逸出物分析对得上，
最后才检查峰位类结论是否有多速率支持。四项中任一缺失，
相应结论都应当降级为「现象描述」。

\srcnote{本节实例分别取自上册 §3.4.1（书内 p.182--189 / PDF p.201--208）、
§3.4.4（书内 p.194--198 / PDF p.213--217）。联用技术对照表与实验条件
影响表中的接口参数、量级范围为学科通用认识整理，源 OCR 未给明文；
正文中所有 \tocheck{} 项待与原页面图像逐条核对。}
\end{crosscheck}
